% GENERATED FILE. Edit canonical lessons, then run npm run generate:books.
% canonical-source-hash: fnv1a64:cf7629b7156a07d9
% canonical-lessons: FA-C37-dadan, FA-C37-bastan, FA-C37-pokhtan, FA-C37-shostan, FA-C37-khandidan

\chapter{To give, To close, To cook, To wash, To laugh}
\label{ch:fa-to-give-to-close-to-cook-to-wash-to-laugh}
\hlchaptermodality{37}

\begin{chapteropening}
\textbf{By the end of this chapter:} \emph{I can say to give, to close, to cook, to wash and to laugh.}

Complete the last lesson of chapter 37: I can say to give, to close, to cook, to wash and to laugh.
\end{chapteropening}

\section[\texorpdfstring{\fa{دادن} (dâdan)}{دادن (dâdan)}]{\fa{دادن} (dâdan) — to give}

\label{lesson:FA-C37-dadan}



\begin{quote}
Before the new one: say the Persian for to buy, then the Persian for to sell.
\end{quote}

\subsection*{You'll want to know: \fa{دادن}}
\textbf{\fa{دادن}} — \emph{dâdan} — \textquotedblleft{}to give\textquotedblright{}.

\begin{cousinweb}
\emph{dâdan} is a native Persian verb, inherited rather than borrowed; it is usually traced to Middle Persian \emph{dādan}. EtymDB takes it back to the Indo-European root \emph{*deh\textsubscript{3}-}, \textquotedblleft{}to give\textquotedblright{}, the root of Latin \emph{dare}.
\end{cousinweb}

\begin{grammarlens}[title={What you've built}]
One more word for this chapter.
\end{grammarlens}

\subsection*{Your turn}
\begin{itemize}
\raggedright
  \item \emph{Say it:} \emph{dâdan}
  \item \emph{Say it:} \emph{dâdan}, once more
  \item \emph{Say it:} say \emph{dâdan}
  \item \emph{Recall:} say the Persian for to buy, then the Persian for to sell, then say \emph{dâdan} again
\end{itemize}

\begin{tcolorbox}[breakable,colback=teal!4,colframe=teal!35!black,title={Before you move on}]
What does \fa{دادن} mean? (\textquotedblleft{}To give\textquotedblright{}.) Say it once more.
\end{tcolorbox}
\section[\texorpdfstring{\fa{بستن} (bastan)}{بستن (bastan)}]{\fa{بستن} (bastan) — to close}

\label{lesson:FA-C37-bastan}



\begin{quote}
Before the new one: say the Persian for to sell, then the Persian for to give.
\end{quote}

\subsection*{You'll want to know: \fa{بستن}}
\textbf{\fa{بستن}} — \emph{bastan} — \textquotedblleft{}to close\textquotedblright{}.

\begin{cousinweb}
\emph{bastan} is a native Persian word, inherited rather than borrowed; it is usually traced to Middle Persian \emph{bastan}, \textquotedblleft{}to bind, to close\textquotedblright{}. EtymDB takes it back to the Indo-European root \emph{*b\textsuperscript{h}end\textsuperscript{h}-}, \textquotedblleft{}to bind\textquotedblright{}, the root of English \emph{bind}.
\end{cousinweb}

\begin{grammarlens}[title={What you've built}]
One more word for this chapter.
\end{grammarlens}

\subsection*{Your turn}
\begin{itemize}
\raggedright
  \item \emph{Say it:} \emph{bastan}
  \item \emph{Say it:} \emph{bastan}, once more
  \item \emph{Say it:} say \emph{bastan}
  \item \emph{Recall:} say the Persian for to sell, then the Persian for to give, then say \emph{bastan} again
\end{itemize}

\begin{tcolorbox}[breakable,colback=teal!4,colframe=teal!35!black,title={Before you move on}]
What does \fa{بستن} mean? (\textquotedblleft{}To close\textquotedblright{}.) Say it once more.
\end{tcolorbox}
\section[\texorpdfstring{\fa{پختن} (pokhtan)}{پختن (pokhtan)}]{\fa{پختن} (pokhtan) — to cook}

\label{lesson:FA-C37-pokhtan}



\begin{quote}
Before the new one: say the Persian for to give, then the Persian for to close.
\end{quote}

\subsection*{You'll want to know: \fa{پختن}}
\textbf{\fa{پختن}} — \emph{pokhtan} — \textquotedblleft{}to cook\textquotedblright{}.

\begin{cousinweb}
\emph{pokhtan} is a native Persian word, inherited rather than borrowed; it is usually traced to Middle Persian \emph{puxtan}, \textquotedblleft{}to cook\textquotedblright{}. EtymDB takes it back to the Indo-European root \emph{*pek\textsuperscript{w}-}, \textquotedblleft{}to cook, to ripen\textquotedblright{}, the root of Latin \emph{coquere}.
\end{cousinweb}

\begin{grammarlens}[title={What you've built}]
One more word for this chapter.
\end{grammarlens}

\subsection*{Your turn}
\begin{itemize}
\raggedright
  \item \emph{Say it:} \emph{pokhtan}
  \item \emph{Say it:} \emph{pokhtan}, once more
  \item \emph{Say it:} say \emph{pokhtan}
  \item \emph{Recall:} say the Persian for to give, then the Persian for to close, then say \emph{pokhtan} again
\end{itemize}

\begin{tcolorbox}[breakable,colback=teal!4,colframe=teal!35!black,title={Before you move on}]
What does \fa{پختن} mean? (\textquotedblleft{}To cook\textquotedblright{}.) Say it once more.
\end{tcolorbox}
\section[\texorpdfstring{\fa{شستن} (shostan)}{شستن (shostan)}]{\fa{شستن} (shostan) — to wash}

\label{lesson:FA-C37-shostan}



\begin{quote}
Before the new one: say the Persian for to close, then the Persian for to cook.
\end{quote}

\subsection*{You'll want to know: \fa{شستن}}
\textbf{\fa{شستن}} — \emph{shostan} — \textquotedblleft{}to wash\textquotedblright{}.

\begin{cousinweb}
\emph{shostan} is a native Persian word, inherited rather than borrowed; it is usually traced to Middle Persian \emph{šustan}, \textquotedblleft{}to wash\textquotedblright{}.
\end{cousinweb}

\begin{grammarlens}[title={What you've built}]
One more word for this chapter.
\end{grammarlens}

\subsection*{Your turn}
\begin{itemize}
\raggedright
  \item \emph{Say it:} \emph{shostan}
  \item \emph{Say it:} \emph{shostan}, once more
  \item \emph{Say it:} say \emph{shostan}
  \item \emph{Recall:} say the Persian for to close, then the Persian for to cook, then say \emph{shostan} again
\end{itemize}

\begin{tcolorbox}[breakable,colback=teal!4,colframe=teal!35!black,title={Before you move on}]
What does \fa{شستن} mean? (\textquotedblleft{}To wash\textquotedblright{}.) Say it once more.
\end{tcolorbox}
\section[\texorpdfstring{\fa{خندیدن} (khandidan)}{خندیدن (khandidan)}]{\fa{خندیدن} (khandidan) — to laugh}

\label{lesson:FA-C37-khandidan}



\begin{quote}
Before the new one: say the Persian for to cook, then the Persian for to wash.
\end{quote}

\subsection*{You'll want to know: \fa{خندیدن}}
\textbf{\fa{خندیدن}} — \emph{khandidan} — \textquotedblleft{}to laugh\textquotedblright{}.

\begin{cousinweb}
\emph{khandidan} is a native Persian word, inherited rather than borrowed; it is usually traced to Middle Persian \emph{xandīdan}, \textquotedblleft{}to laugh\textquotedblright{}.
\end{cousinweb}

\begin{grammarlens}[title={What you've built}]
One more word for this chapter.
\end{grammarlens}

\subsection*{Your turn}
\begin{itemize}
\raggedright
  \item \emph{Say it:} \emph{khandidan}
  \item \emph{Say it:} \emph{khandidan}, once more
  \item \emph{Say it:} say \emph{khandidan}
  \item \emph{Recall:} say the Persian for to cook, then the Persian for to wash, then say \emph{khandidan} again
\end{itemize}

\begin{tcolorbox}[breakable,colback=teal!4,colframe=teal!35!black,title={Before you move on}]
What does \fa{خندیدن} mean? (\textquotedblleft{}To laugh\textquotedblright{}.) Say it once more.
\end{tcolorbox}
